\documentclass[10pt,a4paper]{amsart}
\usepackage[margin=19mm]{geometry}
\usepackage{amsmath,amssymb,mathtools}
\usepackage[scr=rsfs,cal=euler]{mathalfa}
\usepackage{xfrac}
\usepackage[unicode,hidelinks,bookmarks=false]{hyperref}
\newcommand{\ie}{i.e.\ }
\newcommand{\cf}{cf.\ }
\newcommand{\resp}{respectively\ }
\newcommand{\Subsection}{\subsection}
\providecommand{\nobreakdash}{\nobreak\hbox{-}}
\setlength{\emergencystretch}{2em}
\multlinegap0pt
\usepackage{amssymb}
\usepackage{mathtools}
\usepackage[shortlabels]{enumitem}
\newtheorem{lemma}{Lemma}
\title{Latin Squares with Few Transversals}
\author{Zur Luria}
\date{Source: arXiv:2609.08624v1 (CC BY 4.0)}
\begin{document}
\maketitle

\noindent\emph{Source and licence.} Latin Squares with Few Transversals. Version arXiv:2609.08624v1 (CC BY 4.0), distributed by arXiv under the Creative Commons Attribution 4.0 International licence, http://creativecommons.org/licenses/by/4.0/, as recorded in the arXiv licence metadata for this exact version and shown on the abstract page https://arxiv.org/abs/2609.08624v1. Full licence notice: \texttt{licenses/arxiv-2609.08624-CC-BY-4.0.txt}.\par\smallskip

\noindent\textit{Editorial scope.} Counted unit: the article's lemma lem:profile (source0.tex lines 155--165). The article's complete original proof of it is delivered with it (source0.tex lines 167--187, one \texttt{proof} environment of 21 lines). No mathematics is written, repaired or re-derived: the statement and the proof are the article's own lines, in the article's own notation, and no retained line is substituted or re-flowed.  Retained uncounted as the source-local context the proof needs: lines 149--153.  Nothing was omitted from any retained span, so the package carries no omission record.  No retained line contains a citation, so the wrapper carries no bibliography and no bibliography entry.  No retained line contains a figure environment, an image-inclusion command or drawing code, so no image is delivered and the package declares no asset dependency.  Editorial formatting only, in addition: an AMS \texttt{amsart} preamble that restates the article's own declarations and macros used by the retained lines, namely the environments \texttt{lemma} on separate counters, together with the standard packages those lines need.  The retained environments print in this document as Lemma 1 in order of appearance, whereas the article prints the same items with the numbering of its own full document.  The complete native source of the article as distributed in the arXiv e-print for this exact version is bundled unchanged as \texttt{sources/209731/source0.tex}.
\begin{equation}\label{eq:line-sums}
 \sum_jx_{ij}=q,\qquad
 \sum_ix_{ij}=q,\qquad
 \sum_{i+j=s}x_{ij}=q.
\end{equation}

\begin{lemma}[The block profile]\label{lem:profile}
There are nonnegative integers $A,B,C$ with $A+B+C=q$ such that
\begin{equation}\label{eq:profile}
X=
\begin{pmatrix}
 A&C&B\\
 B&A&C\\
 C&B&A
\end{pmatrix}.
\end{equation}
\end{lemma}

\begin{proof}
Set
\[
A=x_{00},\qquad C=x_{01},\qquad B=x_{02}.
\]
The first block-row equation gives $A+B+C=q$. The matrix on the right-hand side of~\eqref{eq:profile} satisfies all the equations in~\eqref{eq:line-sums}, so it remains only to show uniqueness.

Subtract this matrix from $X$. Since both matrices satisfy the row, column and block-symbol sum equations in~\eqref{eq:line-sums}, their difference has sum zero on every block row, block column and block-symbol class. As its first row is zero and its column sums are zero, it has the form
\[
\begin{pmatrix}
0&0&0\\
p&r&s\\
-p&-r&-s
\end{pmatrix}.
\]
Now consider the three block-symbol classes, that is, the sets of positions with $i+j$ fixed modulo $3$. Their sums are zero, and hence
\[
s-r=0,\qquad p-s=0,\qquad r-p=0.
\]
Thus $p=r=s$. The second row sum is also zero, so $3p=0$, and therefore $p=r=s=0$.
\end{proof}

\end{document}
